\documentclass[11pt,a4paper]{article}
\usepackage[utf8]{inputenc}
\usepackage[T1]{fontenc}
\usepackage{amsmath,amssymb,amsthm}
\usepackage{algorithm}
\usepackage{algpseudocode}
\usepackage{geometry}
\usepackage{hyperref}
\usepackage{booktabs}
\usepackage{microtype}

\geometry{margin=1in}

\title{Hierarchical Stability of GRA Obnullification: A Variational Formulation}
\author{bitsoev oleg}
\date{May 2026}

\newtheorem{theorem}{Theorem}
\newtheorem{lemma}[theorem]{Lemma}
\newtheorem{axiom}{Axiom}
\newtheorem{definition}{Definition}

\begin{document}

\maketitle

\begin{abstract}
A formalization of the GRA obnullification mechanism is proposed as a variational hierarchical functional that determines the dynamics of ``foam'' $\Phi$ in a multi-level system. The derivative $d\Phi/dh$ with respect to the level parameter $h$ makes it possible to distinguish productive collapse of contradictions, stagnation, and degradation. Axioms, definitions, lemmas, and a theorem on the stability of obnullification are formulated, as well as a discrete algorithm with pseudocode for practical implementation.
\end{abstract}

\section{Introduction}

The GRA (Gradient Reduction of Argumentative Foam) framework aims to minimize ``foam'' $\Phi$ --- a measure of contradictions, noise, and redundancy in hierarchical systems. The basic obnullification mechanism assumes that the top level of the hierarchy tends to $\Phi \approx 0$, which is interpreted as global coherence. In this work, obnullification is considered not only statically but also dynamically --- through the derivative $d\Phi/dh$. This allows introducing a stability criterion for obnullification and investigating the system's sensitivity to changes in the hierarchy.

\section{Hierarchical Model and Foam Functional}

A system with levels $l = 0,\dots,L$ is considered. The state of a level is $x_l \in \mathbb{R}^{n_l}$. The top level $L$ sets global coherence. Local foam: $\Phi_l = \Phi(x_l) \geq 0$. 

\textbf{Hierarchical GRA-functional:}
\begin{equation}
J[x_0,\dots,x_L] = \sum_{l=0}^L \alpha_l \Phi(x_l) + \sum_{l=0}^{L-1} \beta_l C_l(x_l, A_l x_{l+1}) + \gamma \Psi(x_L),
\label{eq:gra-functional}
\end{equation}
where:
\begin{itemize}
    \item $C_l$ --- penalty for mismatch between adjacent levels,
    \item $A_l$ --- lifting/projection operator,
    \item $\Psi(x_L)$ --- top-level obnullification condition,
    \item $\alpha_l, \beta_l, \gamma > 0$ --- weights.
\end{itemize}

Under continuous parametrization of the level $h$, $\Phi(h)$ is differentiable, and $d\Phi/dh$ is the zero-gradient.

\section{Axioms of GRA Obnullification}

\begin{axiom}[Hierarchy]
The system is organized into levels $l = 0,\dots,L$, each with its own state.
\end{axiom}

\begin{axiom}[Foam]
For each level a foam functional $\Phi(x_l) \geq 0$ is defined.
\end{axiom}

\begin{axiom}[Obnullification]
The operation $T_l$ reduces foam: $\Phi(T_l(x_l)) \leq \Phi(x_l)$.
\end{axiom}

\begin{axiom}[Monotonicity]
In the ideal regime $\Phi_{l+1} \leq \Phi_l$.
\end{axiom}

\begin{axiom}[Sensitivity]
$\Phi(h)$ is differentiable, the sign of $d\Phi/dh$ determines the evolution mode.
\end{axiom}

\section{Definitions}

\begin{definition}[GRA-functional]
See Equation~\ref{eq:gra-functional}.
\end{definition}

\begin{definition}[Stable obnullification]
At a point $h^*$: $\Phi(h^*) = 0$, $\frac{d\Phi}{dh}(h^*) = 0$, $\frac{d^2\Phi}{dh^2}(h^*) > 0$.
\end{definition}

\begin{definition}[Discrete foam derivative]
$\Delta\Phi_l = \Phi_{l+1} - \Phi_l$, its discrete analogue of $d\Phi/dh$.
\end{definition}

\begin{definition}[Zero-gradient]
$\nabla_h \Phi := \frac{d\Phi}{dh}$, characterizes the direction and rate of foam change along the hierarchy.
\end{definition}

\section{Lemmas}

\begin{lemma}[Local stationarity]
At a minimum, the variational derivatives with respect to each $x_l$ vanish.
\end{lemma}

\begin{lemma}[Zero foam does not guarantee stability]
It is also necessary that $\frac{d\Phi}{dh} = 0$.
\end{lemma}

\begin{lemma}[Positive derivative --- degradation]
If $\frac{d\Phi}{dh} > 0$, foam grows.
\end{lemma}

\begin{lemma}[Discrete stability criterion]
For a level $l^*$: $\Phi_{l^*} = 0$, $\Phi_{l^*+1} - \Phi_{l^*} = 0$, $\Phi_{l^*+1} - 2\Phi_{l^*} + \Phi_{l^*-1} > 0$.
\end{lemma}

\section{Stability Theorem}

\begin{theorem}
Let $\Phi(h)$ be twice differentiable with respect to $h$ and let there exist a point $h^*$ such that
\begin{equation}
\Phi(h^*) = 0, \quad \frac{d\Phi}{dh}(h^*) = 0, \quad \frac{d^2\Phi}{dh^2}(h^*) > 0.
\label{eq:stability-conditions}
\end{equation}
Then obnullification at level $h^*$ is stable: in some neighborhood perturbations do not lead to foam growth, and the system returns to $\Phi \approx 0$.
\end{theorem}

\section{Proof Sketch}

\begin{enumerate}
    \item From $\frac{d\Phi}{dh}(h^*) = 0$ it follows that $h^*$ is a critical point.
    \item The condition $\frac{d^2\Phi}{dh^2}(h^*) > 0$ means a local minimum.
    \item There exists a neighborhood $U(h^*)$ where $\Phi(h) \geq \Phi(h^*) = 0$.
    \item Since $\Phi \geq 0$, zero foam is locally minimal.
    \item Under small perturbations a gradient process returns the system to $h^*$, which means stability. In the discrete case the reasoning is repeated via finite differences.
\end{enumerate}

\section{Discrete Algorithm and Pseudocode}

The algorithm combines GRA obnullification, gradient alignment $d\Phi/dh$, and control of the second derivative.

\begin{algorithm}[H]
\caption{HierarchicalStabilityOptimizer}
\begin{algorithmic}[1]
\State \textbf{class} HierarchicalStabilityOptimizer
\State \quad \textbf{def} \_\_init\_\_(self, hierarchy, lr=0.01):
\State \quad \quad self.h = hierarchy \Comment{list of levels with .state and .foam()}
\State \quad \quad self.lr = lr
\State \quad \textbf{def} step(self):
\State \quad \quad \# 1. GRA obnullification
\State \quad \quad \textbf{for} level \textbf{in} self.h: level.state = gra\_nullify(level.state)
\State \quad \quad \# 2. Align dPhi/dh
\State \quad \quad \textbf{for} l \textbf{in} range(len(self.h)-1):
\State \quad \quad \quad dPhi = self.h[l+1].foam() - self.h[l].foam()
\State \quad \quad \quad \textbf{if} abs(dPhi) > 1e-6:
\State \quad \quad \quad \quad grad = grad\_connection(self.h[l], self.h[l+1], dPhi)
\State \quad \quad \quad \quad self.h[l+1].state -= self.lr * grad
\State \quad \quad \# 3. Ensure d\^2Phi/dh\^2 > 0
\State \quad \quad \textbf{for} l \textbf{in} range(1, len(self.h)-1):
\State \quad \quad \quad d2Phi = self.h[l+1].foam() - 2*self.h[l].foam() + self.h[l-1].foam()
\State \quad \quad \quad \textbf{if} d2Phi <= 0:
\State \quad \quad \quad \quad self.h[l].state += self.lr * 0.01 * (self.h[l+1].state + self.h[l-1].state - 2*self.h[l].state)
\State \quad \textbf{def} optimize(self, max\_iter=100):
\State \quad \quad \textbf{for} \_ \textbf{in} range(max\_iter): self.step()
\end{algorithmic}
\end{algorithm}

The full implementation with abstract functions is located in \texttt{src/optimizer.py}.

\section{Discussion and Applications}

Interpretation of the derivative sign:
\begin{itemize}
    \item $\frac{d\Phi}{dh} < 0$ --- productive foam collapse
    \item $\frac{d\Phi}{dh} \approx 0$ --- stagnation or plateau
    \item $\frac{d\Phi}{dh} > 0$ --- degradation (hallucinations, growth of contradictions)
\end{itemize}

\textbf{Applications:}
\begin{itemize}
    \item Diagnostics of cognitive stability of AI architectures (monitoring ``foam'' in hidden states of neural networks)
    \item Multi-level management in organizations (analyzing consistency of decisions across hierarchy levels)
    \item Analysis of the evolution of scientific theories (tracking contradictions in knowledge systems)
\end{itemize}

\section{Installation and Usage}

\begin{verbatim}
git clone https://github.com/qqewq/GRA-Hierarchical-Stability.git
cd GRA-Hierarchical-Stability
pip install -r requirements.txt
python examples/demo.py
\end{verbatim}

To integrate with GRA-Multiverse-Final, replace the stubs in \texttt{src/foam\_utils.py} with real functions from your project.

\section{References}

\begin{itemize}
    \item GRA-Multiverse-Final
    \item GRA-Swarm
    \item NeurIPS 2022 hierarchical optimization materials
    \item HAL preprint hal-04088837
    \item SPP1962 preprint 036, WIAS Berlin
\end{itemize}

\end{document}